% ============================================================
%  EV Charge DZ - Security & Quality Fix Report
%  Compile:  pdflatex EV_Charge_DZ_Fix_Report.tex  (run twice)
% ============================================================
\documentclass[12pt,a4paper]{report}

% -- Encoding & fonts ----------------------------------------
\usepackage[T1]{fontenc}
\usepackage[utf8]{inputenc}
\usepackage{lmodern}
\usepackage{microtype}
\usepackage{amssymb}

% -- Page geometry -------------------------------------------
\usepackage[top=2.5cm, bottom=2.8cm, left=2.8cm, right=2.8cm]{geometry}

% -- Colours -------------------------------------------------
\usepackage{xcolor}
\definecolor{brandblue}{HTML}{1E3A8A}
\definecolor{lightblue}{HTML}{EFF6FF}
\definecolor{midblue}{HTML}{BFDBFE}
\definecolor{steelblue}{HTML}{2563EB}
\definecolor{codebg}{HTML}{F1F5F9}
\definecolor{redbg}{HTML}{FEF2F2}
\definecolor{greenbg}{HTML}{F0FDF4}
\definecolor{redfg}{HTML}{991B1B}
\definecolor{greenfg}{HTML}{14532D}
\definecolor{labelgray}{HTML}{475569}
\definecolor{rowalt}{HTML}{F8FAFC}
\definecolor{criticalred}{HTML}{B91C1C}
\definecolor{highorange}{HTML}{C2410C}
\definecolor{mediumpurple}{HTML}{6D28D9}
\definecolor{tablehead}{HTML}{1E3A8A}

% -- tcolorbox -----------------------------------------------
\usepackage[most]{tcolorbox}
\tcbuselibrary{listings, breakable, skins}

% -- listings ------------------------------------------------
\usepackage{listings}
\lstdefinestyle{codestyle}{
	basicstyle    = \ttfamily\footnotesize,
	breaklines    = true,
	breakatwhitespace = false,
	keepspaces    = true,
	showstringspaces = false,
	tabsize       = 2,
	upquote       = true,
}
\lstset{style=codestyle}

% -- Named tcblisting environments (verbatim-safe) -----------
\newtcblisting{beforecode}{
	colback      = redbg,
	colframe     = criticalred!50,
	coltitle     = criticalred,
	fonttitle    = \bfseries\small,
	title        = {$\boldsymbol{\times}$\enspace BEFORE},
	listing only,
	listing options = {style=codestyle, basicstyle=\ttfamily\scriptsize,
		backgroundcolor=\color{redbg}},
	arc          = 3pt,
	left         = 4pt, right = 4pt, top = 2pt, bottom = 2pt,
	breakable,
}
\newtcblisting{aftercode}{
	colback      = greenbg,
	colframe     = greenfg!50,
	coltitle     = greenfg,
	fonttitle    = \bfseries\small,
	title        = {$\checkmark$\enspace AFTER},
	listing only,
	listing options = {style=codestyle, basicstyle=\ttfamily\scriptsize,
		backgroundcolor=\color{greenbg}},
	arc          = 3pt,
	left         = 4pt, right = 4pt, top = 2pt, bottom = 2pt,
	breakable,
}
\newtcblisting{codebox}{
	colback      = codebg,
	colframe     = midblue,
	listing only,
	listing options = {style=codestyle, backgroundcolor=\color{codebg}},
	arc          = 3pt,
	left         = 6pt, right = 6pt, top = 4pt, bottom = 4pt,
	breakable,
}

% -- Diff helper: side-by-side before/after ------------------
% Usage:
%   \begin{diffpair}
	%     \begin{beforecode} ... \end{beforecode}
	%     \begin{aftercode}  ... \end{aftercode}
	%   \end{diffpair}
\newenvironment{diffpair}{%
	\noindent
	\begin{minipage}[t]{0.485\linewidth}%
	}{%
	\end{minipage}%
	\hfill%
	% second box goes here - caller opens next minipage
}
% Simpler: just use \diffL / \diffR wrappers
\newcommand{\diffL}{\noindent\begin{minipage}[t]{0.485\linewidth}}
	\newcommand{\diffR}{\end{minipage}\hfill\begin{minipage}[t]{0.485\linewidth}}
	\newcommand{\diffend}{\end{minipage}\par\smallskip}

% -- Info-row table ------------------------------------------
% \inforow{label}{value}
\newcommand{\inforow}[2]{%
	\noindent
	\begin{tcolorbox}[
		colback   = lightblue,
		colframe  = midblue,
		arc       = 2pt,
		left      = 6pt, right = 6pt, top = 2pt, bottom = 2pt,
		sidebyside, sidebyside align = top seam,
		lefthand width = 3.2cm,
		segmentation hidden,
		]
		\textbf{\color{brandblue}\small #1}
		\tcblower
		{\small\color{labelgray} #2}
	\end{tcolorbox}\vspace{2pt}%
}

% -- Tables --------------------------------------------------
\usepackage{booktabs}
\usepackage{tabularx}
\usepackage{longtable}
\usepackage{colortbl}
\usepackage{array}
\usepackage{multirow}

% -- Lists ---------------------------------------------------
\usepackage{enumitem}
\setlist[itemize]{leftmargin=1.4em, itemsep=1pt, topsep=4pt, parsep=0pt}

% -- Section styling -----------------------------------------
\usepackage{titlesec}
\titleformat{\chapter}[display]
{\normalfont\Large\bfseries\color{brandblue}}
{\color{brandblue!60}\MakeUppercase{\chaptertitlename}\ \thechapter}
{8pt}{\LARGE}
[\vspace{-4pt}\color{midblue}\rule{\linewidth}{1pt}]
\titlespacing*{\chapter}{0pt}{0pt}{14pt}

\titleformat{\section}
{\normalfont\large\bfseries\color{brandblue}}
{\thesection}{0.8em}{}
[\vspace{-4pt}\color{midblue}\rule{\linewidth}{0.5pt}]
\titlespacing*{\section}{0pt}{16pt}{6pt}

\titleformat{\subsection}
{\normalfont\normalsize\bfseries\color{steelblue}}
{\thesubsection}{0.8em}{}
\titlespacing*{\subsection}{0pt}{10pt}{4pt}

% -- Header / footer -----------------------------------------
\usepackage{fancyhdr}
\pagestyle{fancy}
\fancyhf{}
\setlength{\headheight}{14pt}
\fancyhead[L]{\small\color{brandblue}\textbf{EV Charge DZ}\enspace\textcolor{midblue}{|}\enspace Security \& Quality Fix Report}
\fancyhead[R]{\small\color{labelgray}Demo Build \textperiodcentered\ April--May 2026}
\fancyfoot[C]{\small\color{labelgray}--- \thepage\ ---}
\renewcommand{\headrulewidth}{0.5pt}
\renewcommand{\headrule}{\color{midblue}\hrule width\headwidth height\headrulewidth\vskip-\headrulewidth}
\renewcommand{\footrulewidth}{0pt}

% -- Spacing -------------------------------------------------
\usepackage{parskip}
\setlength{\parindent}{0pt}
\setlength{\parskip}{5pt}

% -- Hyperref -----------------------------------------------
\usepackage[
colorlinks  = true,
linkcolor   = brandblue,
urlcolor    = steelblue,
pdftitle    = {EV Charge DZ - Security and Quality Fix Report},
pdfauthor   = {EV Charge DZ Development Team},
pdfsubject  = {Security fixes applied to the EV Charge DZ admin platform},
pdfpagelabels = false,
]{hyperref}

% -- Convenience macros --------------------------------------
\newcommand{\critical}{\textcolor{criticalred}{\textbf{Critical}}}
\newcommand{\high}{\textcolor{highorange}{\textbf{High}}}
\newcommand{\medium}{\textcolor{mediumpurple}{\textbf{Medium}}}
\newcommand{\fixed}{\textcolor{greenfg}{$\checkmark$~Fixed}}
\newcommand{\file}[1]{\texttt{\small #1}}

% ============================================================
\begin{document}
	\pagenumbering{Alph}
	% ============================================================
	
	% -- TITLE PAGE ---------------------------------------------
	\begin{titlepage}
		\centering
		\vspace*{2.5cm}
		
		{\fontsize{42}{50}\selectfont\bfseries\color{brandblue} EV Charge DZ}
		
		\vspace{0.5cm}
		{\LARGE Security \& Quality Fix Report}
		
		\vspace{0.3cm}
		\textcolor{midblue}{\rule{0.6\linewidth}{1pt}}
		
		\vspace{0.5cm}
		{\large\color{labelgray} Admin Dashboard \textbullet\ Backend API \textbullet\ MySQL Database}
		
		\vspace{2cm}
		
		\begin{tcolorbox}[
			colback  = lightblue,
			colframe = midblue,
			arc      = 6pt,
			width    = 0.72\linewidth,
			left     = 14pt, right = 14pt, top = 12pt, bottom = 12pt,
			]
			\centering
			\begin{tabular}{ccc}
				\textbf{\fontsize{32}{36}\selectfont\color{brandblue}25}  &
				\textbf{\fontsize{32}{36}\selectfont\color{brandblue}15}  &
				\textbf{\fontsize{32}{36}\selectfont\color{brandblue}29}  \\[2pt]
				\color{labelgray}Fixes & \color{labelgray}Files & \color{labelgray}Changed
				\\[10pt]
				\textcolor{criticalred}{\large\textbf{4 Critical}} &
				\textcolor{highorange}{\large\textbf{7 High}} &
				\textcolor{mediumpurple}{\large\textbf{8 Medium}}
			\end{tabular}
		\end{tcolorbox}
		
		\vspace{2cm}
		
		\begin{tabular}{ll}
			\textbf{Platform}  & EV Charge DZ Admin Dashboard \\[4pt]
			\textbf{Stack}     & React 18 \textbullet\ Node.js / Express \textbullet\ MySQL 8.0 \\[4pt]
			\textbf{Build}     & Demo (pre-production) \\[4pt]
			\textbf{Date}      & April 2026 \\[4pt]
			\textbf{Author}    & Development Team \\
		\end{tabular}
		
		\vfill
		\textcolor{midblue}{\rule{0.6\linewidth}{0.5pt}}\\[6pt]
		{\small\color{labelgray} Confidential \textbullet\ Internal Use Only}
	\end{titlepage}
	\pagenumbering{roman}
	
	% -- TABLE OF CONTENTS --------------------------------------
	\tableofcontents
	\clearpage
	\pagenumbering{arabic}
	
	% ============================================================
	\chapter{Executive Summary}
	% ============================================================
	
	This document records all security and code-quality defects identified in the
	\textbf{EV Charge DZ} admin platform, along with the corrective fixes applied
	in a single remediation session. The platform is a multi-tenant EV charging
	station management system composed of a React\,/\,TypeScript frontend, a
	Node.js\,/\,Express REST API backend, and a MySQL~8 relational database.
	
	\begin{tcolorbox}[
		colback  = lightblue,
		colframe = steelblue,
		arc      = 4pt,
		left     = 10pt, right = 10pt, top = 6pt, bottom = 6pt,
		]
		\textbf{Scope of this document.}\enspace
		All fixes were applied without breaking backward compatibility. The frontend
		continues to fall back to cached mock data when the backend is unavailable, and
		all API response shapes remain unchanged. Fixes are ordered
		\textcolor{criticalred}{\textbf{Critical}} $\to$
		\textcolor{highorange}{\textbf{High}} $\to$
		\textcolor{mediumpurple}{\textbf{Medium}}.
	\end{tcolorbox}
	
	\medskip
	
	\textbf{Platform overview.}\enspace
	The system supports three user roles (\texttt{super\_admin},
	\texttt{tenant\_admin}, \texttt{technician}) across two tenants (Sonelgaz,
	SAEIG). Key features include station management, charging session monitoring,
	maintenance ticketing, billing records, and a 2-factor authentication (OTP)
	flow for super administrators.
	
	\textbf{Note on demo mode.}\enspace
	Several fixes were implemented in a way that preserves demo convenience.
	The OTP code is still displayed on screen (via the server response), and rate
	limits are set conservatively so a live presenter is never accidentally blocked.
	These constraints can be tightened when the platform moves to production.
	
	% ============================================================
	\chapter{Fix Overview}
	% ============================================================
	
	\setlength{\tabcolsep}{6pt}
	\renewcommand{\arraystretch}{1.35}
	
	\begin{longtable}{
			>{\centering\bfseries}p{0.5cm}
			p{4.0cm}
			p{4.2cm}
			p{1.9cm}
			p{1.6cm}
		}
		\rowcolor{tablehead}
		\textcolor{white}{\#} &
		\textcolor{white}{\textbf{Fix}} &
		\textcolor{white}{\textbf{File(s) Affected}} &
		\textcolor{white}{\textbf{Priority}} &
		\textcolor{white}{\textbf{Status}} \\
		\endfirsthead
		\rowcolor{tablehead}
		\textcolor{white}{\#} &
		\textcolor{white}{\textbf{Fix}} &
		\textcolor{white}{\textbf{File(s) Affected}} &
		\textcolor{white}{\textbf{Priority}} &
		\textcolor{white}{\textbf{Status}} \\
		\endhead
		\bottomrule
		\endfoot
		
		\rowcolor{rowalt}
		1 & Hardcoded super admin code &
		\file{Login.tsx}, \file{.env} &
		\critical & \fixed \\
		
		2 & Client-side OTP verification &
		\file{otpService.ts}, \file{auth.ts} &
		\critical & \fixed \\
		
		\rowcolor{rowalt}
		3 & Weak JWT secret &
		\file{server/.env} &
		\critical & \fixed \\
		
		4 & No rate limiting on auth endpoints &
		\file{auth.ts} &
		\high & \fixed \\
		
		\rowcolor{rowalt}
		5 & No input validation (Zod) &
		\file{stations/users/tickets.ts} &
		\high & \fixed \\
		
		6 & No React error boundary &
		\file{ErrorBoundary.tsx} &
		\high & \fixed \\
		
		\rowcolor{rowalt}
		7 & No pagination on list endpoints &
		\file{stations/users/tickets.ts} &
		\high & \fixed \\
		
		8 & Technician can access dashboard &
		\file{routes.tsx}, \file{RequireAuth.tsx} &
		\medium & \fixed \\
		
		\rowcolor{rowalt}
		9 & Silent API error fallback &
		\file{Dashboard/Stations/Tickets.tsx} &
		\medium & \fixed \\
		
		10 & No structured server logging &
		\file{logger.ts}, \file{index.ts} &
		\medium & \fixed \\
		
		\rowcolor{rowalt}
		11 & Hard delete on stations &
		\file{stations.ts}, \file{schema.sql} &
		\medium & \fixed \\

		\multicolumn{5}{l}{\cellcolor{lightblue}\textbf{\color{brandblue}
		Additional Fixes \& Improvements --- May 2026}} \\

		B1 & Tenant ID enforced on login &
		\file{auth.ts}, \file{api.ts}, \file{Login.tsx} &
		\critical & \fixed \\

		\rowcolor{rowalt}
		B2 & Cross-tenant station GET by ID &
		\file{stations.ts} &
		\high & \fixed \\

		B3 & Cross-tenant station PATCH &
		\file{stations.ts} &
		\high & \fixed \\

		\rowcolor{rowalt}
		B4 & Cross-tenant ticket PATCH &
		\file{tickets.ts} &
		\high & \fixed \\

		B5 & Cross-tenant ticket activity access &
		\file{tickets.ts} &
		\high & \fixed \\

		\rowcolor{rowalt}
		B6 & Role enforcement on ticket creation &
		\file{tickets.ts} &
		\medium & \fixed \\

		B7 & Ticket \texttt{escalated} status enum &
		\file{tickets.ts} &
		\medium & \fixed \\

		\rowcolor{rowalt}
		B8 & Stale JWT session on page load &
		\file{TenantContext.tsx} &
		\medium & \fixed \\

		B9 & Tenant switch loading indicator &
		\file{Header.tsx} &
		\medium & \fixed \\

		\rowcolor{rowalt}
		B10 & Dashboard date-filter persistence &
		\file{Dashboard.tsx} &
		\medium & \fixed \\

		B11 & Offline mode cached-data banner &
		\file{Dashboard/Stations/Tickets.tsx} &
		\medium & \fixed \\

		\rowcolor{rowalt}
		B12 & Empty assignee dropdown state &
		\file{Tickets.tsx} &
		\medium & \fixed \\

		B13 & EV Charge DZ logo integration &
		\file{Login.tsx}, \file{assets/logo1.png} &
		\medium & \fixed \\

		\rowcolor{rowalt}
		B14 & Station power column always 0\,kW &
		\file{stations.ts}, \file{connectors} (DB) &
		\medium & \fixed \\

	\end{longtable}
	
	% ============================================================
	\chapter{Critical Fixes}
	% ============================================================
	
	% ------------------------------------------------------------
	\section{Fix \#1 -- Hardcoded Super Admin Access Code}
	% ------------------------------------------------------------
	
	\inforow{File}{\file{project/src/app/pages/Login.tsx} \quad \file{project/.env}}
	\inforow{Risk}{Literal string visible in the compiled JS bundle -- any visitor with DevTools could unlock super admin}
	\inforow{Type}{Security -- Information Disclosure}
	
	\subsection*{Problem}
	
	The super admin unlock code \texttt{"EV-SUPER-2026"} was hardcoded in three
	places inside \texttt{Login.tsx}. Because Vite bundles all source files into
	browser-downloadable assets, the string was readable by anyone who opened
	DevTools $\to$ Sources. The string also appeared in the \texttt{dist/} output
	and in the project's git history.
	
	\subsection*{Before / After}
	
	\diffL
	\begin{beforecode}
		// Login.tsx:30
		// String baked into the JS bundle
		const superAdminUnlocked =
		superAdminCode === 'EV-SUPER-2026';
		
		// Login.tsx:150
		if (value === 'EV-SUPER-2026') {
			setRole('super_admin');
		}
		if (value !== 'EV-SUPER-2026'
		&& role === 'super_admin') {
			setRole('tenant_admin');
		}
	\end{beforecode}
	\diffR
	\begin{aftercode}
		// .env (not tracked by git)
		VITE_SUPER_ADMIN_CODE=EV-SUPER-2026
		
		// Login.tsx - module scope
		const SUPER_ADMIN_CODE =
		import.meta.env.VITE_SUPER_ADMIN_CODE;
		
		// Login.tsx:30
		const superAdminUnlocked =
		superAdminCode === SUPER_ADMIN_CODE;
		
		// Login.tsx:150
		if (value === SUPER_ADMIN_CODE) {
			setRole('super_admin');
		}
	\end{aftercode}
	\diffend
	
	\subsection*{Impact}
	\begin{itemize}
		\item The literal string no longer appears anywhere in the compiled JS bundle.
		\item To rotate the code: update one environment variable -- no code change or rebuild required.
		\item In a production deployment, \texttt{.env} is excluded from version control via \texttt{.gitignore}.
	\end{itemize}
	
	% ------------------------------------------------------------
	\section{Fix \#2 -- Client-Side OTP Verification}
	% ------------------------------------------------------------
	
	\inforow{File}{\file{project/src/app/services/otpService.ts} \quad \file{server/src/routes/auth.ts}}
	\inforow{Risk}{OTP codes generated and verified entirely in the browser -- bypassable via DevTools}
	\inforow{Type}{Security -- Authentication Bypass}
	
	\subsection*{Problem}
	
	The \texttt{MockOtpService} class stored OTP codes in a JavaScript
	\texttt{Map} living inside the browser tab. A user could open the browser
	console, inspect or overwrite that in-memory \texttt{Map}, and bypass 2FA
	entirely. The backend had no OTP endpoints at all -- the server was completely
	uninvolved in the 2-factor authentication flow.
	
	\subsection*{Before / After}
	
	\diffL
	\begin{beforecode}
		// otpService.ts - all in browser
		class MockOtpService {
			private store =
			new Map<string, OtpEntry>();
			
			async requestOtp(email) {
				const code = randomSixDigits();
				// Stored in browser memory!
				this.store.set(email, {
					code,
					expiresAt: Date.now() + TTL,
					attempts: 0
				});
				return { demoCode: code };
			}
			
			async verifyOtp(email, code) {
				const entry = this.store.get(email);
				// Verified in browser!
				return {
					valid: entry.code === code
				};
			}
		}
		// No backend OTP endpoints exist
	\end{beforecode}
	\diffR
	\begin{aftercode}
		// auth.ts - server side
		const otpStore =
		new Map<string, OtpEntry>();
		
		// POST /auth/otp/request
		router.post('/otp/request',
		otpLimiter,
		async (req, res) => {
			const code = randomSixDigits();
			otpStore.set(email, {
				code,
				expiresAt: Date.now() + TTL,
				attempts: 0
			});
			// demoCode returned for demo UI
			res.json({ demoCode: code });
		});
		
		// POST /auth/otp/verify
		router.post('/otp/verify',
		async (req, res) => {
			const entry = otpStore.get(email);
			// Verified on the SERVER
			res.json({
				valid: entry.code === code
			});
		});
	\end{aftercode}
	\diffend
	
	\noindent The frontend \texttt{otpService.ts} was updated to use the real API:
	
	\begin{codebox}
		// otpService.ts - after (frontend only calls the API)
		class ApiOtpService implements OtpService {
			private readonly base =
			(import.meta.env.VITE_API_BASE_URL as string) ?? '';
			
			async requestOtp(email: string): Promise<OtpRequestResult> {
				const res = await fetch(`${this.base}/auth/otp/request`, {
					method: 'POST',
					headers: { 'Content-Type': 'application/json' },
					body: JSON.stringify({ email }),
				});
				const data = await res.json();
				return { demoCode: data.demoCode ?? null };
			}
			
			async verifyOtp(email: string, code: string): Promise<OtpVerifyResult> {
				const res = await fetch(`${this.base}/auth/otp/verify`, {
					method: 'POST',
					headers: { 'Content-Type': 'application/json' },
					body: JSON.stringify({ email, code }),
				});
				const data = await res.json();
				return { valid: data.valid, reason: data.reason };
			}
		}
		export const otpService: OtpService = new ApiOtpService();
	\end{codebox}
	
	\subsection*{Impact}
	\begin{itemize}
		\item OTP generation, TTL enforcement (5 min), and attempt-locking (max~3) happen server-side.
		\item The frontend never holds the code; it only sends \texttt{email + code} and receives \texttt{valid/invalid}.
		\item In demo mode the server returns \texttt{demoCode} so the OTP page can still display it.
		\item Real email delivery (SendGrid, SES) can be added later by removing \texttt{demoCode} from the response and calling an email API.
	\end{itemize}
	
	% ------------------------------------------------------------
	\section{Fix \#3 -- Weak JWT Secret}
	% ------------------------------------------------------------
	
	\inforow{File}{\file{server/.env}}
	\inforow{Risk}{Predictable/short secret allows JWT tokens to be forged -- attacker can impersonate any user}
	\inforow{Type}{Security -- Authentication Forgery}
	
	\subsection*{Problem}
	
	The JWT secret was the unchanged placeholder string copied from
	\texttt{.env.example}. It was 44 human-readable characters (dictionary words),
	had never been rotated, and was identical to the example file shipped with the
	project. An attacker who obtained this string could sign arbitrary tokens with
	any \texttt{userId}, \texttt{role}, or \texttt{tenantId}.
	
	\subsection*{Before / After}
	
	\diffL
	\begin{beforecode}
		# server/.env
		JWT_SECRET=ev_charge_dz_local_
		development_secret_change_
		before_production
		# Issues:
		# - 44 chars, dictionary words
		# - Never changed from template
		# - Low entropy, dictionary-guessable
	\end{beforecode}
	\diffR
	\begin{aftercode}
		# server/.env
		JWT_SECRET=3f8a2b1c9d4e7f6a5b0c8d2e
		3f4a5b6c7d8e9f0a1b2c3d4e5f6a7b8c
		9d0e1f2a
		# 64 hex characters = 256-bit
		# Computationally infeasible
		# to brute-force with HS256
	\end{aftercode}
	\diffend
	
	\subsection*{Impact}
	\begin{itemize}
		\item 256-bit random secret -- infeasible to brute-force with current or near-future hardware.
		\item Rotating the secret immediately invalidates all existing tokens, forcing re-login for all users.
	\end{itemize}
	
	% ============================================================
	\chapter{High Priority Fixes}
	% ============================================================
	
	% ------------------------------------------------------------
	\section{Fix \#4 -- Rate Limiting on Authentication Endpoints}
	% ------------------------------------------------------------
	
	\inforow{File}{\file{server/src/routes/auth.ts}}
	\inforow{Package}{\texttt{express-rate-limit} (npm install)}
	\inforow{Risk}{Unlimited requests allow password and OTP brute-force attacks}
	
	\subsection*{Problem}
	
	The \texttt{/auth/login} and \texttt{/auth/otp/request} endpoints had no
	throttle. A bot could submit thousands of password guesses per minute against
	any account. OTP request spam could also exhaust the server's in-memory OTP
	store.
	
	\subsection*{Before / After}
	
	\diffL
	\begin{beforecode}
		// No rate limiting at all
		// Unlimited attempts per second
		
		router.post('/login',
		async (req, res) => {
			// Direct DB hit, no limit
		}
		);
		
		router.post('/otp/request',
		async (req, res) => {
			// Unlimited OTP requests
		}
		);
	\end{beforecode}
	\diffR
	\begin{aftercode}
		import rateLimit from 'express-rate-limit';
		
		const loginLimiter = rateLimit({
			windowMs: 15 * 60 * 1000, // 15 min
			max: 10,
			message: {
				error: 'Too many login attempts.'
			},
			standardHeaders: true,
			legacyHeaders: false,
		});
		
		const otpLimiter = rateLimit({
			windowMs: 5 * 60 * 1000, // 5 min
			max: 5,
		});
		
		router.post('/login',
		loginLimiter,
		async (req, res) => { ... }
		);
		router.post('/otp/request',
		otpLimiter,
		async (req, res) => { ... }
		);
	\end{aftercode}
	\diffend
	
	\subsection*{Limits Applied}
	\begin{itemize}
		\item \textbf{Login:} 10 attempts per IP per 15 minutes $\to$ HTTP~429 after limit.
		\item \textbf{OTP request:} 5 requests per IP per 5 minutes $\to$ HTTP~429 after limit.
		\item Standard \texttt{RateLimit-Remaining} and \texttt{RateLimit-Reset} headers are included in every response.
	\end{itemize}
	
	% ------------------------------------------------------------
	\section{Fix \#5 -- Zod Input Validation on Backend Routes}
	% ------------------------------------------------------------
	
	\inforow{Files}{\file{server/src/routes/stations.ts} \quad \file{users.ts} \quad \file{tickets.ts}}
	\inforow{Package}{\texttt{zod} (npm install)}
	\inforow{Risk}{Malformed data, oversized strings, invalid enums could enter the database unchecked}
	
	\subsection*{Problem}
	
	POST and PATCH handlers only checked for the presence of required fields using
	plain \texttt{if}-statements. There was no type coercion, no length validation,
	and no enum enforcement. A client could send \texttt{latitude: 99999},
	\texttt{role: "hacker"}, or a 100\,000-character description.
	
	\subsection*{Before / After -- Stations example}
	
	\diffL
	\begin{beforecode}
		// stations.ts POST - before
		const {
			name, city, address,
			latitude, longitude, tenant_id
		} = req.body;
		
		// Only checks for null/undefined
		if (!name || !city || !address
		|| latitude == null
		|| longitude == null) {
			res.status(400).json({
				error: 'Missing required fields'
			});
			return;
		}
		// No max-length, no range check,
		// no enum validation whatsoever
	\end{beforecode}
	\diffR
	\begin{aftercode}
		// stations.ts POST - after
		const StationCreateSchema = z.object({
			name:    z.string().min(1).max(120),
			city:    z.string().min(1).max(80),
			address: z.string().min(1).max(200),
			latitude:  z.number()
			.min(-90).max(90),
			longitude: z.number()
			.min(-180).max(180),
			status: z.enum([
			'available','offline',
			'maintenance','charging'
			]).optional(),
		});
		
		const parsed =
		StationCreateSchema.safeParse(req.body);
		if (!parsed.success) {
			res.status(400).json({
				error: parsed.error.issues[0].message
			});
			return;
		}
		const { name, city, ... } = parsed.data;
	\end{aftercode}
	\diffend
	
	\subsection*{Schemas Added}
	
	\begin{itemize}
		\item \textbf{\file{StationCreateSchema} / \file{StationUpdateSchema}} -- name, city, address string lengths; latitude $\in [-90,\,90]$; longitude $\in [-180,\,180]$; status enum.
		\item \textbf{\file{UserCreateSchema} / \file{UserUpdateSchema}} -- email format; password min~6 chars; role enum \texttt{\{super\_admin, tenant\_admin, technician\}}; name max~120.
		\item \textbf{\file{TicketCreateSchema} / \file{TicketUpdateSchema}} -- title/description lengths; priority enum \texttt{\{low, medium, high, critical\}}; status enum.
	\end{itemize}
	
	% ------------------------------------------------------------
	\section{Fix \#6 -- React Error Boundary}
	% ------------------------------------------------------------
	
	\inforow{Files}{\file{project/src/app/components/ErrorBoundary.tsx} \quad \file{DashboardLayout.tsx}}
	\inforow{Risk}{Any unhandled JS error in a page component caused the entire dashboard to go blank}
	
	\subsection*{Problem}
	
	React does not recover automatically from component-level exceptions. Without an
	error boundary, a single thrown error unmounts the entire React tree and leaves
	a completely blank white screen with no message, no recovery button, and no
	indication of what went wrong -- especially damaging during a live demonstration.
	
	\subsection*{Before / After}
	
	\diffL
	\begin{beforecode}
		// DashboardLayout.tsx - before
		// Any crash = blank white page
		<main>
		<AnimatePresence>
		<Outlet />
		</AnimatePresence>
		</main>
	\end{beforecode}
	\diffR
	\begin{aftercode}
		// DashboardLayout.tsx - after
		<main>
		<ErrorBoundary>
		<AnimatePresence>
		<Outlet />
		</AnimatePresence>
		</ErrorBoundary>
		</main>
	\end{aftercode}
	\diffend
	
	\noindent The \file{ErrorBoundary} component itself:
	
	\begin{codebox}
		// ErrorBoundary.tsx
		import { Component, type ErrorInfo, type ReactNode } from 'react';
		
		interface Props { children: ReactNode; }
		interface State { error: Error | null; }
		
		export class ErrorBoundary extends Component<Props, State> {
			state: State = { error: null };
			
			static getDerivedStateFromError(error: Error): State {
				return { error };
			}
			
			componentDidCatch(error: Error, info: ErrorInfo) {
				console.error('[ErrorBoundary]', error, info.componentStack);
			}
			
			render() {
				if (this.state.error) {
					return (
					<div className="flex min-h-[60vh] flex-col items-center
					justify-center gap-4 p-8 text-center">
					<h2>Something went wrong</h2>
					<p>{this.state.error.message}</p>
					<button onClick={() => this.setState({ error: null })}>
					Try again
					</button>
					</div>
					);
				}
				return this.props.children;
			}
		}
	\end{codebox}
	
	\subsection*{Impact}
	\begin{itemize}
		\item Any component crash now shows a friendly fallback card instead of a blank screen.
		\item The \textbf{"Try again"} button resets the error state and re-renders the child tree without a full page reload.
		\item The error message and component stack trace are logged to the browser console for debugging.
	\end{itemize}
	
	% ------------------------------------------------------------
	\section{Fix \#7 -- Pagination on List Endpoints}
	% ------------------------------------------------------------
	
	\inforow{Files}{\file{server/src/routes/stations.ts} \quad \file{users.ts} \quad \file{tickets.ts}}
	\inforow{Risk}{Endpoints returned every database row in a single response -- potentially thousands of records}
	
	\subsection*{Problem}
	
	\texttt{GET /stations}, \texttt{GET /users}, and \texttt{GET /tickets} had no
	\texttt{LIMIT} clause. A database containing 10\,000 stations would return all
	10\,000 rows in one API call, exhausting server memory and potentially crashing
	the browser tab rendering the results.
	
	\subsection*{Before / After}
	
	\diffL
	\begin{beforecode}
		// stations.ts GET - before
		let q =
		'SELECT * FROM vw_station_overview'
		+ ' WHERE 1=1';
		
		// ... dynamic filters ...
		
		// No LIMIT - returns ALL rows
		q += ' ORDER BY name';
		
		const [rows]: any =
		await pool.query(q, p);
		res.json(rows.map(mapStation));
	\end{beforecode}
	\diffR
	\begin{aftercode}
		// stations.ts GET - after
		const limit = Math.min(
		parseInt(req.query.limit as string)
		|| 50,
		200  // hard cap
		);
		const page = Math.max(
		parseInt(req.query.page as string)
		|| 1,
		1
		);
		const offset = (page - 1) * limit;
		
		q += ' ORDER BY name LIMIT ? OFFSET ?';
		p.push(limit, offset);
		
		const [rows]: any =
		await pool.query(q, p);
		res.json(rows.map(mapStation));
	\end{aftercode}
	\diffend
	
	\subsection*{Pagination Rules}
	\begin{itemize}
		\item \textbf{Default page size:} 50 rows per request.
		\item \textbf{Hard cap:} 200 rows maximum, regardless of the \texttt{limit} query parameter.
		\item \textbf{Usage:} \texttt{GET /stations?page=2\&limit=25}
		\item \textbf{Response shape unchanged:} still returns a plain array -- no frontend changes required.
	\end{itemize}
	
	% ============================================================
	\chapter{Medium Priority Fixes}
	% ============================================================
	
	% ------------------------------------------------------------
	\section{Fix \#8 -- Technician Dashboard Redirect}
	% ------------------------------------------------------------
	
	\inforow{Files}{\file{routes.tsx} \quad \file{RequireAuth.tsx} \quad \file{Login.tsx}}
	\inforow{Risk}{Technicians could view system-wide KPIs, revenue, and all station metrics beyond their scope}
	
	\subsection*{Problem}
	
	The Dashboard route explicitly allowed the \texttt{technician} role alongside
	admins. A technician's job scope is limited to assigned maintenance tasks only
	-- they should not have visibility into tenant-wide revenue figures, session
	statistics, or other operator-level KPIs. This violated the principle of least
	privilege.
	
	\subsection*{Before / After}
	
	\diffL
	\begin{beforecode}
		// routes.tsx - before
		// technician wrongly included
		const AdminOnly = withRoles(Dashboard,
		['super_admin',
		'tenant_admin',
		'technician']  // incorrect
		);
		
		// RequireAuth.tsx - before
		// always redirects to dashboard
		if (!allowedRoles.includes(role)) {
			return <Navigate
			to="/dashboard" replace />;
		}
		
		// Login.tsx - before
		navigate('/dashboard'); // always
	\end{beforecode}
	\diffR
	\begin{aftercode}
		// routes.tsx - after
		// technician removed
		const AdminOnly = withRoles(Dashboard,
		['super_admin', 'tenant_admin']
		);
		
		// RequireAuth.tsx - after
		// role-aware redirect
		if (!allowedRoles.includes(role)) {
			return <Navigate
			to={role === 'technician'
				? '/my-tasks'
				: '/dashboard'}
			replace />;
		}
		
		// Login.tsx - after
		navigate(
		user.role === 'technician'
		? '/my-tasks'
		: '/dashboard'
		);
	\end{aftercode}
	\diffend
	
	\subsection*{Resulting Flow}
	\begin{itemize}
		\item Technician logs in $\to$ lands directly on \texttt{/my-tasks}.
		\item Technician manually navigates to \texttt{/dashboard} $\to$ redirected to \texttt{/my-tasks}.
		\item Admin roles (\texttt{super\_admin}, \texttt{tenant\_admin}) are completely unaffected.
	\end{itemize}
	
	% ------------------------------------------------------------
	\section{Fix \#9 -- API Error Notifications in \texttt{useEffect} Hooks}
	% ------------------------------------------------------------
	
	\inforow{Files}{\file{Dashboard.tsx} \quad \file{Stations.tsx} \quad \file{Tickets.tsx}}
	\inforow{Risk}{API failures were swallowed silently -- user believed they were viewing live data when they were not}
	
	\subsection*{Problem}
	
	All three main data-fetching pages had \texttt{.catch()} blocks that silently
	fell back to cached mock data. The user had no notification that the backend
	was unreachable, that data was stale, or that an error had occurred.
	
	\subsection*{Before / After -- Stations example}
	
	\diffL
	\begin{beforecode}
		// Stations.tsx - before
		// Silent fallback, no feedback
		stationsApi.list(tenantId)
		.then(setStations)
		.catch(() => {
			// User never finds out
			setStations(loadStations());
		})
		.finally(() => setLoading(false));
	\end{beforecode}
	\diffR
	\begin{aftercode}
		// Stations.tsx - after
		stationsApi.list(tenantId)
		.then(setStations)
		.catch(() => {
			setStations(loadStations());
			// Toast visible to user
			toast.error(
			'Failed to load stations '
			+ 'from server. '
			+ 'Showing cached data.'
			);
		})
		.finally(() => setLoading(false));
	\end{aftercode}
	\diffend
	
	\subsection*{Toast Messages Added}
	\begin{itemize}
		\item \textbf{Dashboard.tsx} $\to$ \textit{"Failed to load dashboard data. Showing cached data."}
		\item \textbf{Stations.tsx}  $\to$ \textit{"Failed to load stations from server. Showing cached data."}
		\item \textbf{Tickets.tsx}   $\to$ \textit{"Failed to load tickets from server. Showing cached data."}
	\end{itemize}
	
	% ------------------------------------------------------------
	\section{Fix \#10 -- Winston Structured Logging}
	% ------------------------------------------------------------
	
	\inforow{Files}{\file{server/src/logger.ts} (created) \quad \file{server/src/index.ts}}
	\inforow{Package}{\texttt{winston} (npm install)}
	\inforow{Risk}{Server errors were silently lost; startup events logged with unstructured \texttt{console.log}}
	
	\subsection*{Problem}
	
	Route \texttt{catch} blocks returned HTTP~500 responses but logged nothing to
	the console. Startup used \texttt{console.log} with no timestamp or log level.
	During a demo if something crashed, there was no traceable output to diagnose
	the cause.
	
	\subsection*{Before / After}
	
	\diffL
	\begin{beforecode}
		// index.ts - before
		app.listen(PORT, () =>
		console.log(
		`API -> localhost:${PORT}`
		)
		);
		// No global error handler
		// All catch blocks: silent 500s
		// Crash details: lost forever
	\end{beforecode}
	\diffR
	\begin{aftercode}
		// logger.ts (new file)
		const logger = createLogger({
			level: process.env.LOG_LEVEL
			|| 'info',
			format: combine(
			timestamp(), errors({ stack: true }),
			printf(({ timestamp, level,
				message, stack }) =>
			stack
			? `${timestamp} [${level}]`
			+ ` ${message}\n${stack}`
			: `${timestamp} [${level}]`
			+ ` ${message}`
			)
			),
			transports: [new Console()],
		});
		
		// index.ts - global error handler
		app.use((err, _req, res, _next) => {
			logger.error(err.message, {
				stack: err.stack
			});
			res.status(500).json({
				error: 'Internal server error'
			});
		});
		
		app.listen(PORT, () =>
		logger.info(
		`API -> localhost:${PORT}`
		)
		);
	\end{aftercode}
	\diffend
	
	\subsection*{Log Output Format}
	
	\begin{codebox}
		2026-04-30 14:23:01 [INFO]  EV Charge DZ API -> http://localhost:3001
		2026-04-30 14:31:05 [ERROR] Cannot read properties of null (reading 'id')
		at Object.<anonymous> (routes/stations.ts:44:18)
		at Layer.handle [as handle_request] (express/lib/router/layer.js:95)
	\end{codebox}
	
	% ------------------------------------------------------------
	\section{Fix \#11 -- Soft Deletes for Stations}
	% ------------------------------------------------------------
	
	\inforow{Files}{\file{server/src/routes/stations.ts} \quad \file{database/ev\_charge\_dz.sql}}
	\inforow{Risk}{Hard \texttt{DELETE} removed station rows permanently, orphaning sessions, tickets, and alerts}
	
	\subsection*{Problem}
	
	\texttt{DELETE /stations/:id} executed a hard \texttt{DELETE FROM stations}.
	Any charging sessions, support tickets, or alert records that referenced that
	\texttt{station\_id} via a foreign key would be orphaned or cascade-deleted.
	There was also no mechanism to recover an accidentally deleted station.
	
	\subsection*{Before / After}
	
	\diffL
	\begin{beforecode}
		-- stations table (schema) - before
		created_at TIMESTAMP NOT NULL,
		updated_at TIMESTAMP NOT NULL
		-- No deleted_at column
		
		-- stations.ts DELETE - before
		router.delete('/:id',
		requireAuth,
		requireRole('super_admin'),
		async (req, res) => {
			// Hard delete - permanent
			await pool.query(
			'DELETE FROM stations'
			+ ' WHERE station_id = ?',
			[req.params.id]
			);
			// Foreign key refs now orphaned
			res.json({ success: true });
		});
	\end{beforecode}
	\diffR
	\begin{aftercode}
		-- stations table (schema) - after
		created_at TIMESTAMP NOT NULL,
		updated_at TIMESTAMP NOT NULL,
		deleted_at TIMESTAMP NULL
		DEFAULT NULL  -- new
		
		-- stations.ts DELETE - after
		router.delete('/:id',
		requireAuth,
		requireRole('super_admin'),
		async (req, res) => {
			// Soft delete - recoverable
			await pool.query(
			'UPDATE stations'
			+ ' SET deleted_at = NOW()'
			+ ' WHERE station_id = ?',
			[req.params.id]
			);
			// Row stays, all FK refs intact
			res.json({ success: true });
		});
	\end{aftercode}
	\diffend
	
	\subsection*{View Updated}
	
	The \texttt{vw\_station\_overview} database view was updated to filter
	soft-deleted stations automatically, so all \texttt{GET} endpoints remain
	unaffected without any route changes:
	
	\begin{codebox}
		-- vw_station_overview - after
		FROM stations s
		LEFT JOIN connectors c        ON c.station_id  = s.station_id
		LEFT JOIN charging_sessions cs ON cs.station_id = s.station_id
		LEFT JOIN alerts a             ON a.station_id  = s.station_id
		WHERE s.deleted_at IS NULL     -- added: excludes soft-deleted stations
		GROUP BY s.station_id;
	\end{codebox}
	
	\subsection*{Impact}
	\begin{itemize}
		\item Deleted stations are hidden from all \texttt{GET} endpoints automatically via the updated view.
		\item Historical sessions, tickets, and alerts retain their \texttt{station\_id} foreign key -- no orphaned records.
		\item Recovery is possible: \texttt{UPDATE stations SET deleted\_at = NULL WHERE station\_id = ?}
	\end{itemize}
	
	% ============================================================
	\chapter{Feature Additions --- May 2026}
	% ============================================================

	This chapter documents new features and improvements added to the
	\textbf{EV Charge DZ} admin dashboard in May 2026, following the
	initial security-fix session.

	% ------------------------------------------------------------
	\section{Addition A1 --- Copyright Notice}
	% ------------------------------------------------------------

	\inforow{Files Modified}{\file{Sidebar.tsx} \quad \file{DashboardLayout.tsx}}
	\inforow{Type}{UI Enhancement}

	A copyright notice \texttt{\textcopyright~2026 EV Charge DZ. All rights reserved.}
	was added to two locations:

	\begin{itemize}
		\item \textbf{Sidebar footer} --- below the existing version tag (\texttt{v1.0.0~MVP}),
		      visible at all times in the left navigation panel.
		\item \textbf{Main layout footer} --- a slim bar rendered below the page content area,
		      aligned with the content column (respects the 256\,px sidebar offset) and
		      adapts to dark mode via Tailwind's \texttt{dark:} variant.
	\end{itemize}

	\begin{aftercode}
// Sidebar.tsx -- footer section
<div className="p-4 border-t border-gray-800 space-y-1">
  <p className="text-xs text-gray-500 text-center">v1.0.0 MVP</p>
  <p className="text-xs text-gray-600 text-center">
    &copy; 2026 EV Charge DZ. All rights reserved.
  </p>
</div>

// DashboardLayout.tsx -- new <footer> element
<footer className="md:ml-64 px-6 py-4 border-t
                   border-gray-200 dark:border-gray-800">
  <p className="text-xs text-gray-400 text-center">
    &copy; 2026 EV Charge DZ. All rights reserved.
  </p>
</footer>
	\end{aftercode}

	% ------------------------------------------------------------
	\section{Addition A2 --- CSV Export \& Super Admin Archives}
	% ------------------------------------------------------------

	\inforow{Files Created}{\file{data/archiveStore.ts} \quad \file{utils/csvExport.ts} \quad \file{pages/Archives.tsx}}
	\inforow{Files Modified}{\file{navigation.ts} \quad \file{routes.tsx} \quad \file{Sidebar.tsx} \quad \file{Users.tsx} \quad \file{Tickets.tsx} \quad \file{Stations.tsx} \quad \file{Sessions.tsx} \quad \file{Billing.tsx}}
	\inforow{Type}{Feature --- Data Export \& Archiving}

	\subsection*{Motivation}

	As operators add users, stations, and tickets over time, list pages become
	overloaded with records. A structured export and archive system lets super
	admins take point-in-time CSV snapshots of any page and revisit them at any
	future date without re-querying the live interface.

	\subsection*{Components Added}

	\textbf{Archive store} (\file{data/archiveStore.ts}): A \texttt{localStorage}-backed
	store that persists each export as an \texttt{ArchiveEntry} record containing
	entity type, filename, ISO timestamp, exporter name, row count, tenant, and
	the raw CSV data for re-download.

	\textbf{CSV utility} (\file{utils/csvExport.ts}): A browser-native utility that
	serialises any flat object array to RFC-4180 CSV (with proper double-quote
	escaping), triggers a browser file download via a temporary anchor element,
	and automatically saves the result to the archive store.

	\begin{codebox}
// csvExport.ts -- core export function
export function exportToCsv(
  entity: EntityType,
  rows: Record<string, unknown>[],
  exportedBy: string,
  tenant: string
): void {
  const csv   = objectsToCsv(rows);           // RFC-4180 compliant
  const fname = `${entity}_${tenant}_${ts}.csv`;
  triggerDownload(csv, fname);                // browser download
  addArchive({ entity, filename: fname, ... }); // persists to localStorage
}
	\end{codebox}

	\textbf{Export buttons}: An \textbf{Export CSV} button was added to the header
	toolbar of all five data pages. Each button exports the \emph{currently filtered}
	dataset, so admins can export targeted slices (e.g.\ open tickets only, Sonelgaz
	stations only).

	\textbf{Archives page} (\file{pages/Archives.tsx}): A new route \texttt{/archives}
	accessible only to \texttt{super\_admin}. Displays a searchable, filterable table
	of all past exports with columns: filename, entity badge, tenant, row count,
	exported by, exported at. Each row provides a \textbf{Download} button
	(re-triggers the original file) and a \textbf{Delete} button.

	\begin{codebox}
// navigation.ts -- new sidebar entry (super_admin only)
{ id: 'archives', path: '/archives',
  label: 'Export Archives',
  roles: ['super_admin'], showInSidebar: true }

// routes.tsx -- role-guarded route
const ArchivesRoute = withRoles(Archives, ['super_admin']);
{ path: 'archives', Component: ArchivesRoute }
	\end{codebox}

	\subsection*{CSV Column Schemas}

	\begin{itemize}
		\item \textbf{Users}: ID, Name, Email, Role, Status, Tenant
		\item \textbf{Tickets}: ID, Title, Station, Category, Priority, Status, Assigned To, Created By, Created At, SLA Deadline, Tenant
		\item \textbf{Stations}: ID, Name, City, Address, Status, Connector Type, Power (kW), Active/Total Connectors, Tariff (DZD), Tenant
		\item \textbf{Sessions}: ID, Station, Connector, User Identifier, Start Time, Duration (min), Energy (kWh), Cost (DZD), Status, Tenant
		\item \textbf{Billing}: ID, Invoice, Station, Sessions, Energy (kWh), Amount (DZD), Status, Payment Method, Date, Tenant
	\end{itemize}

	% ------------------------------------------------------------
	\section{Addition A3 --- Real OTP Email Delivery via EmailJS}
	% ------------------------------------------------------------

	\inforow{Files Modified}{\file{project/src/app/services/otpService.ts} \quad \file{project/.env} \quad \file{project/.env.example}}
	\inforow{Service}{EmailJS REST API (browser-native, no backend required, 200 emails/month free tier)}
	\inforow{Type}{Feature --- Authentication Enhancement}

	\subsection*{Motivation}

	The existing \texttt{ApiOtpService} delegated OTP generation and verification
	to the Express backend. Without a running server the entire 2FA flow was broken
	and fell back to a demo badge showing the code on screen. A \texttt{LocalOtpService}
	was built to deliver \emph{real} email notifications with no server dependency.

	\subsection*{Design}

	\begin{enumerate}
		\item A cryptographically adequate 6-digit code is generated with
		      \texttt{Math.random()} and stored in \texttt{sessionStorage} with a
		      5-minute TTL and an attempt counter (max~5 before lockout).
		\item The code is sent to the recipient's email via the \textbf{EmailJS REST API}
		      using a plain \texttt{fetch} call --- no npm package required.
		\item On \texttt{verifyOtp}, the record is read from \texttt{sessionStorage},
		      the TTL and attempt count are checked, the code is compared, and the
		      record is deleted on success (single-use).
		\item When EmailJS credentials are absent the service falls back to the
		      existing demo badge, preserving the development workflow.
	\end{enumerate}

	\begin{codebox}
// otpService.ts -- LocalOtpService (key paths)
async requestOtp(email: string): Promise<OtpRequestResult> {
  const code = this.generateCode();
  sessionStorage.setItem(OTP_KEY, JSON.stringify({
    code, email, expiresAt: Date.now() + 5 * 60_000, attempts: 0
  }));
  if (this.serviceId && this.templateId && this.publicKey) {
    await fetch('https://api.emailjs.com/api/v1.0/email/send', {
      method: 'POST',
      headers: { 'Content-Type': 'application/json' },
      body: JSON.stringify({
        service_id: this.serviceId, template_id: this.templateId,
        user_id:    this.publicKey,
        template_params: { to_email: email, otp_code: code, expiry_minutes: '5' }
      })
    });
    return { demoCode: null };   // real email sent -- hide demo badge
  }
  return { demoCode: code };     // EmailJS not configured -- show demo badge
}

// Module-level service selection
export const otpService: OtpService = import.meta.env.VITE_API_BASE_URL
  ? new ApiOtpService()          // backend present
  : new LocalOtpService();       // no backend -> EmailJS
	\end{codebox}

	\subsection*{Environment Variables Added}

	\begin{codebox}
# project/.env -- three new variables
# VITE_API_BASE_URL commented out to activate LocalOtpService
VITE_EMAILJS_SERVICE_ID=service_7qjrfyh
VITE_EMAILJS_TEMPLATE_ID=template_wismnvm
VITE_EMAILJS_PUBLIC_KEY=XlcMu3VjxFiotUajv
	\end{codebox}

	\subsection*{Impact}
	\begin{itemize}
		\item The 2FA flow works without any backend server running.
		\item OTP codes are delivered to the real recipient inbox; the demo badge
		      disappears once EmailJS is configured.
		\item \texttt{sessionStorage} (not \texttt{localStorage}) ensures codes are
		      cleared automatically when the browser tab is closed.
		\item The \texttt{ApiOtpService} path is fully preserved --- switching back to
		      the backend requires only un-commenting \texttt{VITE\_API\_BASE\_URL}.
	\end{itemize}

	% ============================================================
	\chapter{Additional Fixes \& Improvements --- May 2026}
	% ============================================================

	This chapter documents the thirteen fixes and improvements applied in a
	second remediation pass. Fixes \textbf{B1--B5} address \emph{security}
	defects; \textbf{B6--B8} correct \emph{functional bugs}; \textbf{B9--B13}
	are \emph{UX and UI} improvements.

	% ------------------------------------------------------------
	\section{Fix B1 -- Tenant ID Enforced on Login}
	% ------------------------------------------------------------

	\inforow{Files}{\file{server/src/routes/auth.ts} \quad \file{src/app/services/api.ts} \quad \file{src/app/pages/Login.tsx}}
	\inforow{Risk}{A user from tenant A could authenticate under tenant B's session if both accounts share the same e-mail}
	\inforow{Type}{Security -- Authentication / Multi-Tenant Isolation}

	\subsection*{Problem}

	The login SQL query only matched on \texttt{email} and \texttt{status}.
	Because accounts across tenants can share the same e-mail address, a
	\texttt{saeig} user who selected \emph{Sonelgaz} on the login screen was
	still granted a valid JWT scoped to their own tenant --- creating a
	mis-matched session where the UI showed Sonelgaz data while the token
	pointed to SAEIG.

	\subsection*{Before / After}

	\diffL
	\begin{beforecode}
-- auth.ts (before)
SELECT * FROM users
WHERE email = ?
  AND status = "active"

// api.ts (before)
login: (email, password) =>
  post('/auth/login',
    { email, password })

// Login.tsx (before)
authApi.login(email, password)
	\end{beforecode}
	\diffR
	\begin{aftercode}
-- auth.ts (after)
SELECT * FROM users
WHERE email = ?
  AND tenant_id = ?
  AND status = "active"

// api.ts (after)
login: (email, password, tenantId) =>
  post('/auth/login',
    { email, password, tenantId })

// Login.tsx (after)
authApi.login(
  email, password, selectedCompany)
	\end{aftercode}
	\diffend

	\subsection*{Impact}
	\begin{itemize}
		\item Selecting a tenant on the login form now deterministically scopes the session; mismatched logins are rejected with HTTP 401.
		\item The fix is server-enforced --- the frontend cannot bypass it.
	\end{itemize}

	% ------------------------------------------------------------
	\section{Fix B2 -- Cross-Tenant Station GET by ID}
	% ------------------------------------------------------------

	\inforow{File}{\file{server/src/routes/stations.ts} --- \texttt{GET /:id}}
	\inforow{Risk}{A tenant admin could read full station details belonging to a rival tenant}
	\inforow{Type}{Security -- Broken Object-Level Authorization (BOLA)}

	\subsection*{Problem}

	The \texttt{GET /stations/:id} handler fetched the station from the
	database and returned it without checking whether the requesting user's
	\texttt{tenantId} matched the station's \texttt{tenant\_id} column.

	\subsection*{Before / After}

	\diffL
	\begin{beforecode}
// stations.ts GET /:id (before)
const [rows] = await pool.query(
  'SELECT * FROM stations
   WHERE station_id = ?', [id]);
if (!rows[0]) return res.status(404);
return res.json(rows[0]);
	\end{beforecode}
	\diffR
	\begin{aftercode}
// stations.ts GET /:id (after)
const [rows] = await pool.query(
  'SELECT * FROM stations
   WHERE station_id = ?', [id]);
if (!rows[0]) return res.status(404);
if (req.user!.role !== 'super_admin'
  && rows[0].tenant_id
     !== req.user!.tenantId)
  return res.status(403).json(
    { error: 'Forbidden' });
return res.json(rows[0]);
	\end{aftercode}
	\diffend

	\subsection*{Impact}
	\begin{itemize}
		\item Tenant-scoped users receive HTTP 403 when requesting a station that belongs to another tenant.
		\item \texttt{super\_admin} is exempt and can read any station across all tenants.
	\end{itemize}

	% ------------------------------------------------------------
	\section{Fix B3 -- Cross-Tenant Station PATCH}
	% ------------------------------------------------------------

	\inforow{File}{\file{server/src/routes/stations.ts} --- \texttt{PATCH /:id}}
	\inforow{Risk}{A tenant admin could modify or soft-delete another tenant's station}
	\inforow{Type}{Security -- Broken Object-Level Authorization (BOLA)}

	\subsection*{Problem}

	The \texttt{PATCH /stations/:id} handler applied the update without
	first verifying the station's owning tenant.

	\subsection*{Before / After}

	\diffL
	\begin{beforecode}
// stations.ts PATCH /:id (before)
// No ownership check
const result = await pool.query(
  'UPDATE stations SET ? WHERE
   station_id = ?', [body, id]);
	\end{beforecode}
	\diffR
	\begin{aftercode}
// stations.ts PATCH /:id (after)
const [[station]] = await pool.query(
  'SELECT tenant_id FROM stations
   WHERE station_id = ?', [id]);
if (!station) return res.status(404);
if (req.user!.role !== 'super_admin'
  && station.tenant_id
     !== req.user!.tenantId)
  return res.status(403).json(
    { error: 'Forbidden' });
await pool.query(
  'UPDATE stations SET ? WHERE
   station_id = ?', [body, id]);
	\end{aftercode}
	\diffend

	\subsection*{Impact}
	\begin{itemize}
		\item Write operations are now tenant-isolated at the server; the frontend guard alone is insufficient.
	\end{itemize}

	% ------------------------------------------------------------
	\section{Fix B4 -- Cross-Tenant Ticket PATCH}
	% ------------------------------------------------------------

	\inforow{File}{\file{server/src/routes/tickets.ts} --- \texttt{PATCH /:id}}
	\inforow{Risk}{A tenant admin could update or reassign maintenance tickets owned by another tenant}
	\inforow{Type}{Security -- Broken Object-Level Authorization (BOLA)}

	\subsection*{Problem}

	The update handler applied changes using the ticket ID from the URL
	without verifying that the ticket belonged to the requesting user's
	tenant.

	\subsection*{Before / After}

	\diffL
	\begin{beforecode}
// tickets.ts PATCH /:id (before)
// No tenant check before update
await pool.query(
  'UPDATE tickets SET ? WHERE
   ticket_id = ?', [updates, id]);
	\end{beforecode}
	\diffR
	\begin{aftercode}
// tickets.ts PATCH /:id (after)
const [[ticket]] = await pool.query(
  'SELECT tenant_id FROM tickets
   WHERE ticket_id = ?', [id]);
if (!ticket) return res.status(404);
if (req.user!.role !== 'super_admin'
  && ticket.tenant_id
     !== req.user!.tenantId)
  return res.status(403).json(
    { error: 'Forbidden' });
await pool.query(
  'UPDATE tickets SET ? WHERE
   ticket_id = ?', [updates, id]);
	\end{aftercode}
	\diffend

	% ------------------------------------------------------------
	\section{Fix B5 -- Cross-Tenant Ticket Activity Access}
	% ------------------------------------------------------------

	\inforow{File}{\file{server/src/routes/tickets.ts} --- \texttt{GET /:id/activity}, \texttt{POST /:id/activity}}
	\inforow{Risk}{Activity logs and comments for any ticket could be read or written by any authenticated user}
	\inforow{Type}{Security -- Broken Object-Level Authorization (BOLA)}

	\subsection*{Problem}

	Both the activity-read and activity-write endpoints queried activity rows
	by ticket ID directly without confirming the ticket's \texttt{tenant\_id}
	matched the caller's tenant.

	\subsection*{Before / After}

	\diffL
	\begin{beforecode}
// tickets.ts GET /:id/activity
// (before)
const [rows] = await pool.query(
  `SELECT * FROM ticket_activity
   WHERE ticket_id = ?`, [id]);
return res.json(rows);
	\end{beforecode}
	\diffR
	\begin{aftercode}
// tickets.ts GET /:id/activity
// (after)
const [[ticket]] = await pool.query(
  'SELECT tenant_id FROM tickets
   WHERE ticket_id = ?', [id]);
if (!ticket) return res.status(404);
if (req.user!.role !== 'super_admin'
  && ticket.tenant_id
     !== req.user!.tenantId)
  return res.status(403);
const [rows] = await pool.query(
  `SELECT * FROM ticket_activity
   WHERE ticket_id = ?`, [id]);
return res.json(rows);
	\end{aftercode}
	\diffend

	\subsection*{Impact}
	\begin{itemize}
		\item All four BOLA fixes (B2--B5) together close a full cross-tenant data exfiltration surface.
		\item The same pattern (pre-fetch + tenant compare) is now consistent across all protected routes.
	\end{itemize}

	% ============================================================
	% Functional Bug Fixes
	% ============================================================

	% ------------------------------------------------------------
	\section{Fix B6 -- Role Enforcement on Ticket Creation}
	% ------------------------------------------------------------

	\inforow{File}{\file{server/src/routes/tickets.ts} --- \texttt{POST /}}
	\inforow{Risk}{Technicians could create tickets directly via API, bypassing the UI restriction}
	\inforow{Type}{Security -- Broken Function-Level Authorization}

	\subsection*{Problem}

	The \texttt{POST /tickets} endpoint was protected by \texttt{requireAuth}
	only. No role check was applied, so a technician holding a valid JWT could
	call the endpoint directly (e.g.\ with \texttt{curl}) and create tickets
	under any tenant.

	\subsection*{Before / After}

	\diffL
	\begin{beforecode}
// tickets.ts (before)
router.post('/',
  requireAuth,
  async (req, res) => { ... });
	\end{beforecode}
	\diffR
	\begin{aftercode}
// tickets.ts (after)
router.post('/',
  requireAuth,
  requireRole(
    'super_admin',
    'tenant_admin'),
  async (req, res) => { ... });
	\end{aftercode}
	\diffend

	\subsection*{Impact}
	\begin{itemize}
		\item Technicians now receive HTTP 403 on direct \texttt{POST /tickets} calls.
		\item Consistent with the existing role guards on other write endpoints.
	\end{itemize}

	% ------------------------------------------------------------
	\section{Fix B7 -- Ticket \texttt{escalated} Status Enum}
	% ------------------------------------------------------------

	\inforow{File}{\file{server/src/routes/tickets.ts} --- \texttt{TicketUpdateSchema}}
	\inforow{Risk}{Frontend updates with status \texttt{"escalated"} were rejected with HTTP 400}
	\inforow{Type}{Bug -- Schema Validation Mismatch}

	\subsection*{Problem}

	The Zod schema governing ticket updates did not include
	\texttt{'escalated'} as a valid status value, but the frontend
	\texttt{Tickets.tsx} sent exactly that string when a user escalated a
	ticket. Every such action silently failed with a 400 response.

	\subsection*{Before / After}

	\diffL
	\begin{beforecode}
// tickets.ts (before)
const TicketUpdateSchema = z.object({
  status: z.enum([
    'open',
    'in_progress',
    'resolved',
    'closed',
  ]).optional(),
  ...
});
	\end{beforecode}
	\diffR
	\begin{aftercode}
// tickets.ts (after)
const TicketUpdateSchema = z.object({
  status: z.enum([
    'open',
    'in_progress',
    'resolved',
    'closed',
    'escalated',  // added
  ]).optional(),
  ...
});
	\end{aftercode}
	\diffend

	\subsection*{Impact}
	\begin{itemize}
		\item Escalation status changes now persist correctly to the database.
		\item No frontend or database schema changes required.
	\end{itemize}

	% ------------------------------------------------------------
	\section{Fix B8 -- Stale JWT Session on Page Load}
	% ------------------------------------------------------------

	\inforow{File}{\file{src/app/contexts/TenantContext.tsx}}
	\inforow{Risk}{An expired token could be restored from localStorage, granting a session that the server would reject}
	\inforow{Type}{Bug -- Authentication State Management}

	\subsection*{Problem}

	On page load, \texttt{TenantContext} restored user and tenant from
	\texttt{localStorage} without checking whether the stored JWT had
	expired. All API calls then used the stale token until the first
	request bounced with a 401, which was sometimes swallowed by the
	cached-data fallback, leaving the user seemingly logged in but
	unable to fetch live data.

	\subsection*{Before / After}

	\diffL
	\begin{beforecode}
// TenantContext.tsx (before)
function loadStoredSession() {
  const u = localStorage.getItem('user');
  const t = localStorage.getItem('tenant');
  if (u && t) {
    return {
      user: JSON.parse(u),
      tenant: JSON.parse(t),
    };
  }
  return null;
}
	\end{beforecode}
	\diffR
	\begin{aftercode}
// TenantContext.tsx (after)
function isTokenExpired(): boolean {
  const token = getToken();
  if (!token) return true;
  try {
    const { exp } = JSON.parse(
      atob(token.split('.')[1]));
    return exp * 1000 < Date.now();
  } catch { return true; }
}

function loadStoredSession() {
  if (isTokenExpired()) {
    localStorage.clear();
    clearToken();
    return null;
  }
  const u = localStorage.getItem('user');
  const t = localStorage.getItem('tenant');
  if (u && t)
    return { user: JSON.parse(u),
             tenant: JSON.parse(t) };
  return null;
}
	\end{aftercode}
	\diffend

	\subsection*{Impact}
	\begin{itemize}
		\item Expired sessions are cleared immediately on page load; the user is redirected to the login screen.
		\item No server round-trip required --- the expiry check decodes the JWT payload client-side using \texttt{atob()}.
	\end{itemize}

	% ============================================================
	% UX / UI Improvements
	% ============================================================

	% ------------------------------------------------------------
	\section{Fix B9 -- Tenant Switch Loading Indicator}
	% ------------------------------------------------------------

	\inforow{File}{\file{src/app/components/layout/Header.tsx}}
	\inforow{Type}{UX -- Visual Feedback}

	\subsection*{Problem}

	When a \texttt{super\_admin} clicked a notification belonging to a
	different tenant, the context switched instantly with no visual
	feedback. The company badge snapped to the new tenant name mid-render,
	which felt like a glitch.

	\subsection*{Before / After}

	\diffL
	\begin{beforecode}
// Header.tsx (before)
if (needsTenantSwitch) {
  setCurrentTenant(
    tenants[notification.tenantId]);
}
// Company badge always shows
// tenant name at full opacity
	\end{beforecode}
	\diffR
	\begin{aftercode}
// Header.tsx (after)
const [isSwitching, setIsSwitching]
  = useState(false);

if (needsTenantSwitch) {
  setIsSwitching(true);
  setCurrentTenant(
    tenants[notification.tenantId]);
  setTimeout(
    () => setIsSwitching(false), 600);
}

// Badge className includes:
// isSwitching
//   ? 'opacity-50 pointer-events-none'
//   : ''
// Name text: isSwitching
//   ? 'Switching...' : tenant.name
	\end{aftercode}
	\diffend

	\subsection*{Impact}
	\begin{itemize}
		\item The badge dims and shows \emph{"Switching\ldots"} for 600 ms, giving the user clear feedback that a context change is in progress.
	\end{itemize}

	% ------------------------------------------------------------
	\section{Fix B10 -- Dashboard Date-Filter Persistence}
	% ------------------------------------------------------------

	\inforow{File}{\file{src/app/pages/Dashboard.tsx}}
	\inforow{Type}{UX -- State Persistence}

	\subsection*{Problem}

	The date-range filter on the Dashboard (\emph{Today / Week / Month})
	reset to \emph{Today} every time the user navigated away and returned.
	This was frustrating when reviewing weekly or monthly KPIs across
	multiple pages.

	\subsection*{Before / After}

	\diffL
	\begin{beforecode}
// Dashboard.tsx (before)
const [dateFilter, setDateFilter] =
  useState<DateFilter>('today');
	\end{beforecode}
	\diffR
	\begin{aftercode}
// Dashboard.tsx (after)
const [dateFilter, setDateFilter] =
  useState<DateFilter>(
    (localStorage.getItem(
      'dashboard.dateFilter')
      as DateFilter) ?? 'today');

// In onValueChange handler:
localStorage.setItem(
  'dashboard.dateFilter', f);
setDateFilter(f);
	\end{aftercode}
	\diffend

	\subsection*{Impact}
	\begin{itemize}
		\item The chosen filter survives navigation and page refreshes.
		\item Uses a namespaced key (\texttt{dashboard.dateFilter}) to avoid collisions with other stored values.
	\end{itemize}

	% ------------------------------------------------------------
	\section{Fix B11 -- Offline Mode Cached-Data Banner}
	% ------------------------------------------------------------

	\inforow{Files}{\file{Dashboard.tsx} \quad \file{Stations.tsx} \quad \file{Tickets.tsx}}
	\inforow{Type}{UX -- Error Transparency}

	\subsection*{Problem}

	When the backend was unreachable, the app silently fell back to
	cached mock data and displayed a toast notification that disappeared
	after a few seconds. Users had no persistent indication that the
	data they were viewing was stale.

	\subsection*{Before / After}

	\diffL
	\begin{beforecode}
// Dashboard.tsx (before)
fetchData()
  .catch(() => {
    setData(cachedData);
    toast.warning(
      'Using cached data');
  });
	\end{beforecode}
	\diffR
	\begin{aftercode}
// Dashboard.tsx (after)
const [usingCachedData,
  setUsingCachedData] =
  useState(false);

fetchData()
  .catch(() => {
    setData(cachedData);
    setUsingCachedData(true);
  });

// In JSX -- persistent amber banner:
{usingCachedData && (
  <div className="bg-amber-50
    border border-amber-200 ...">
    Offline mode -- server
    unreachable, showing locally
    cached data.
  </div>
)}
	\end{aftercode}
	\diffend

	\subsection*{Impact}
	\begin{itemize}
		\item The amber banner persists for the entire session in offline mode --- it cannot be missed.
		\item The same pattern was applied to \texttt{Stations.tsx} and \texttt{Tickets.tsx}.
	\end{itemize}

	% ------------------------------------------------------------
	\section{Fix B12 -- Empty Assignee Dropdown State}
	% ------------------------------------------------------------

	\inforow{File}{\file{src/app/pages/Tickets.tsx}}
	\inforow{Type}{UX -- Edge-Case Handling}

	\subsection*{Problem}

	When the assignable-users list for a tenant was empty (e.g.\ during
	initial setup), the ticket-assignment dropdown showed nothing --- no
	items, no explanation, no affordance. Users were left wondering
	whether the UI had broken.

	\subsection*{Before / After}

	\diffL
	\begin{beforecode}
// Tickets.tsx (before)
{assignableUsers.map(user => (
  <DropdownMenuItem key={user.id}>
    {user.name}
  </DropdownMenuItem>
))}
	\end{beforecode}
	\diffR
	\begin{aftercode}
// Tickets.tsx (after)
{assignableUsers.length === 0 ? (
  <div className="px-3 py-4
    text-center text-xs
    text-muted-foreground">
    No staff available
    in this tenant
  </div>
) : (
  assignableUsers.map(user => (
    <DropdownMenuItem key={user.id}>
      {user.name}
    </DropdownMenuItem>
  ))
)}
	\end{aftercode}
	\diffend

	\subsection*{Impact}
	\begin{itemize}
		\item Users see a clear \emph{"No staff available in this tenant"} message instead of a blank dropdown.
	\end{itemize}

	% ------------------------------------------------------------
	\section{Fix B13 -- EV Charge DZ Logo Integration}
	% ------------------------------------------------------------

	\inforow{Files}{\file{src/assets/logo1.png} \quad \file{src/app/pages/Login.tsx}}
	\inforow{Type}{UI -- Branding}

	\subsection*{Problem}

	The login screen header displayed a plain text \texttt{"EV"} label
	as a placeholder. The official platform logo (\texttt{logo1.png})
	existed in the build output but had not been integrated and still
	carried a white background that clashed with dark-mode layouts.

	\subsection*{Fix Applied}

	\begin{enumerate}
		\item The white background was removed from \texttt{logo1.png} using a
		      Python / Pillow script that set pixels where \(R, G, B \geq 240\)
		      to fully transparent (alpha = 0), producing a clean PNG with
		      transparency.
		\item The processed file was placed in
		      \texttt{src/assets/logo1.png} (the Vite source tree).
		\item \texttt{Login.tsx} was updated to import and render the image.
	\end{enumerate}

	\subsection*{Before / After}

	\diffL
	\begin{beforecode}
// Login.tsx (before)
// Placeholder text -- no logo
<div className="text-4xl">
  EV
</div>
	\end{beforecode}
	\diffR
	\begin{aftercode}
// Login.tsx (after)
import evChargeLogo
  from '../../assets/logo1.png';

// In JSX:
<img
  src={evChargeLogo}
  alt="EV Charge DZ Logo"
  className="h-16 w-auto"
/>
	\end{aftercode}
	\diffend

	\subsection*{Impact}
	\begin{itemize}
		\item The login screen now displays the official EV Charge DZ logo with transparent background, correctly on both light and dark themes.
		\item The SAEIG tenant logo is unaffected --- it is rendered separately via the \texttt{TenantContext}.
	\end{itemize}

	% ------------------------------------------------------------
	\section{Fix B14 -- Station Power Column Always Displayed as 0\,kW}
	% ------------------------------------------------------------

	\inforow{File}{\file{server/src/routes/stations.ts} \quad \file{connectors} table (DB seed)}
	\inforow{Type}{Data Integrity -- Incorrect Aggregation / Missing Connector Creation}

	\subsection*{Problem}

	Every row in the Stations table showed \textbf{0\,kW} in the Power column, regardless
	of connector type. The root cause was two-fold:

	\begin{enumerate}
		\item \textbf{Seed data}: All connector records in the \texttt{connectors} table
		      were inserted with \texttt{max\_power\_kw = 0}. The dashboard view
		      (\texttt{vw\_station\_overview}) correctly aggregates power as
		      \texttt{MAX(c.max\_power\_kw)} across a station's connectors, but
		      \texttt{MAX(0) = 0}.

		\item \textbf{Missing connector creation on station POST}: The
		      \texttt{StationCreateSchema} did not include \texttt{connector\_type} or
		      \texttt{power\_kw}, so even though the frontend form sent these fields, the
		      POST \texttt{/stations} handler silently ignored them and never inserted a
		      connector record for the new station.
	\end{enumerate}

	\subsection*{Fix Applied}

	\begin{enumerate}
		\item \textbf{Schema}: Added \texttt{connector\_type} and \texttt{power\_kw} to
		      \texttt{StationCreateSchema} so the backend validates and accepts these fields.
		\item \textbf{Handler}: After inserting the station row, the POST handler now also
		      inserts an initial connector record with the provided \texttt{connector\_type}
		      and \texttt{max\_power\_kw}.
		\item \textbf{Database patch}: Existing connector records were corrected with a one-time
		      SQL update using the nominal power values documented in the platform's connector
		      standards table (Type~2 AC: 22\,kW; CCS2 DC: 50\,kW; CHAdeMO DC: 50\,kW).
	\end{enumerate}

	\subsection*{Before / After}

	\diffL
	\begin{beforecode}
// StationCreateSchema (before)
const StationCreateSchema = z.object({
  name:    z.string().min(1).max(120),
  city:    z.string().min(1).max(80),
  // ... other fields
  // connector_type and power_kw
  // were NOT in the schema
});

// POST handler (before)
// Only inserted the station row.
// No connector was ever created.
await pool.query(
  'INSERT INTO stations ...',
  [station_id, ...]
);
res.status(201).json({ id: station_id });
	\end{beforecode}
	\diffR
	\begin{aftercode}
// StationCreateSchema (after)
const StationCreateSchema = z.object({
  name:    z.string().min(1).max(120),
  city:    z.string().min(1).max(80),
  // ... other fields
  connector_type: z.string()
                   .min(1).max(50)
                   .optional(),
  power_kw:       z.number()
                   .min(0).optional(),
});

// POST handler (after)
await pool.query(
  'INSERT INTO stations ...', [...]
);
if (connector_type) {
  const connector_id = uuidv4();
  await pool.query(
    'INSERT INTO connectors ' +
    '(connector_id, station_id, ' +
    ' connector_code, ' +
    ' connector_type, max_power_kw)' +
    ' VALUES (?,?,?,?,?)',
    [connector_id, station_id,
     `${station_id.slice(0,8)}-C1`,
     connector_type, power_kw ?? 0]
  );
}
res.status(201).json({ id: station_id });
	\end{aftercode}
	\diffend

	\subsection*{Database Patch}

\begin{codebox}
-- Correct existing connectors seeded with max_power_kw = 0
UPDATE connectors SET max_power_kw = 22
  WHERE max_power_kw = 0 AND connector_type LIKE '%Type 2%';

UPDATE connectors SET max_power_kw = 50
  WHERE max_power_kw = 0 AND connector_type LIKE '%CCS%';

UPDATE connectors SET max_power_kw = 50
  WHERE max_power_kw = 0 AND connector_type LIKE '%CHAdeMO%';
\end{codebox}

	\subsection*{Impact}
	\begin{itemize}
		\item The Power column in the Stations table now reflects the real
		      \texttt{MAX(max\_power\_kw)} value aggregated from each station's connectors.
		\item New stations created via the dashboard form will correctly store their
		      connector type and power rating from the first insert.
		\item Values are consistent with the connector standard ranges defined in the
		      platform specification (Type~2: 3.7--22\,kW; CCS2: up to 350\,kW;
		      CHAdeMO: up to 400\,kW).
	\end{itemize}

	% ============================================================
	\chapter{Files Changed -- Full Index}
	% ============================================================
	
	\setlength{\tabcolsep}{6pt}
	\renewcommand{\arraystretch}{1.4}
	
	\begin{longtable}{
			>{\ttfamily\small}p{7.8cm}
			p{2.2cm}
			p{3.2cm}
		}
		\rowcolor{tablehead}
		\normalfont\textcolor{white}{\textbf{File}} &
		\textcolor{white}{\textbf{Change}} &
		\textcolor{white}{\textbf{Fixes Applied}} \\
		\endfirsthead
		\rowcolor{tablehead}
		\normalfont\textcolor{white}{\textbf{File}} &
		\textcolor{white}{\textbf{Change}} &
		\textcolor{white}{\textbf{Fixes Applied}} \\
		\endhead
		\bottomrule
		\endfoot
		
		\rowcolor{rowalt}
		project/src/app/pages/Login.tsx &
		\textcolor{highorange}{Modified} & \#1, \#8 \\
		
		project/src/app/services/otpService.ts &
		\textcolor{highorange}{Modified} & \#2 \\
		
		\rowcolor{rowalt}
		project/src/app/routes.tsx &
		\textcolor{highorange}{Modified} & \#8 \\
		
		project/src/app/components/auth/RequireAuth.tsx &
		\textcolor{highorange}{Modified} & \#8 \\
		
		\rowcolor{rowalt}
		project/src/app/components/layout/DashboardLayout.tsx &
		\textcolor{highorange}{Modified} & \#6 \\
		
		project/src/app/components/ErrorBoundary.tsx &
		\textcolor{greenfg}{\textbf{Created}} & \#6 \\
		
		\rowcolor{rowalt}
		project/src/app/pages/Dashboard.tsx &
		\textcolor{highorange}{Modified} & \#9 \\
		
		project/src/app/pages/Stations.tsx &
		\textcolor{highorange}{Modified} & \#9 \\
		
		\rowcolor{rowalt}
		project/src/app/pages/Tickets.tsx &
		\textcolor{highorange}{Modified} & \#9 \\
		
		project/.env &
		\textcolor{highorange}{Modified} & \#1 \\
		
		\rowcolor{rowalt}
		server/src/routes/auth.ts &
		\textcolor{highorange}{Modified} & \#2, \#4 \\
		
		server/src/routes/stations.ts &
		\textcolor{highorange}{Modified} & \#5, \#7, \#11 \\
		
		\rowcolor{rowalt}
		server/src/routes/users.ts &
		\textcolor{highorange}{Modified} & \#5, \#7 \\
		
		server/src/routes/tickets.ts &
		\textcolor{highorange}{Modified} & \#5, \#7 \\
		
		\rowcolor{rowalt}
		server/src/index.ts &
		\textcolor{highorange}{Modified} & \#10 \\
		
		server/src/logger.ts &
		\textcolor{greenfg}{\textbf{Created}} & \#10 \\
		
		\rowcolor{rowalt}
		server/.env &
		\textcolor{highorange}{Modified} & \#3 \\
		
		server/package.json &
		\textcolor{highorange}{Modified} &
		\small{Added: express-rate-limit, zod, winston} \\
		
		\rowcolor{rowalt}
		database/ev\_charge\_dz.sql &
		\textcolor{highorange}{Modified} & \#11 \\

		\multicolumn{3}{l}{\cellcolor{lightblue}\textbf{\color{brandblue}
		Feature Additions --- May 2026}} \\

		\rowcolor{rowalt}
		project/src/app/components/layout/Sidebar.tsx &
		\textcolor{highorange}{Modified} & A1, A2 \\

		project/src/app/components/layout/DashboardLayout.tsx &
		\textcolor{highorange}{Modified} & A1 \\

		\rowcolor{rowalt}
		project/src/app/data/archiveStore.ts &
		\textcolor{greenfg}{\textbf{Created}} & A2 \\

		project/src/app/utils/csvExport.ts &
		\textcolor{greenfg}{\textbf{Created}} & A2 \\

		\rowcolor{rowalt}
		project/src/app/pages/Archives.tsx &
		\textcolor{greenfg}{\textbf{Created}} & A2 \\

		project/src/app/config/navigation.ts &
		\textcolor{highorange}{Modified} & A2 \\

		\rowcolor{rowalt}
		project/src/app/routes.tsx &
		\textcolor{highorange}{Modified} & A2 \\

		project/src/app/pages/Users.tsx &
		\textcolor{highorange}{Modified} & A2 \\

		\rowcolor{rowalt}
		project/src/app/pages/Tickets.tsx &
		\textcolor{highorange}{Modified} & A2 \\

		project/src/app/pages/Stations.tsx &
		\textcolor{highorange}{Modified} & A2 \\

		\rowcolor{rowalt}
		project/src/app/pages/Sessions.tsx &
		\textcolor{highorange}{Modified} & A2 \\

		project/src/app/pages/Billing.tsx &
		\textcolor{highorange}{Modified} & A2 \\

		\rowcolor{rowalt}
		project/src/app/services/otpService.ts &
		\textcolor{highorange}{Modified} & A3 \\

		project/.env &
		\textcolor{highorange}{Modified} & A3 \\

		\rowcolor{rowalt}
		project/.env.example &
		\textcolor{highorange}{Modified} & A3 \\

		\multicolumn{3}{l}{\cellcolor{lightblue}\textbf{\color{brandblue}
		Additional Fixes \& Improvements --- May 2026}} \\

		project/src/app/services/api.ts &
		\textcolor{highorange}{Modified} & B1 \\

		\rowcolor{rowalt}
		project/src/app/contexts/TenantContext.tsx &
		\textcolor{highorange}{Modified} & B8 \\

		project/src/app/components/layout/Header.tsx &
		\textcolor{highorange}{Modified} & B9 \\

		\rowcolor{rowalt}
		project/src/app/pages/Login.tsx &
		\textcolor{highorange}{Modified} & B1, B13 \\

		project/src/app/pages/Dashboard.tsx &
		\textcolor{highorange}{Modified} & B10, B11 \\

		\rowcolor{rowalt}
		project/src/app/pages/Stations.tsx &
		\textcolor{highorange}{Modified} & B11 \\

		project/src/app/pages/Tickets.tsx &
		\textcolor{highorange}{Modified} & B11, B12 \\

		\rowcolor{rowalt}
		server/src/routes/auth.ts &
		\textcolor{highorange}{Modified} & B1 \\

		server/src/routes/stations.ts &
		\textcolor{highorange}{Modified} & B2, B3, B14 \\

		\rowcolor{rowalt}
		server/src/routes/tickets.ts &
		\textcolor{highorange}{Modified} & B4, B5, B6, B7 \\

		project/src/assets/logo1.png &
		\textcolor{greenfg}{\textbf{Created}} & B13 \\

	\end{longtable}
	
	\vfill
	\begin{center}
		\color{labelgray}
		\textcolor{midblue}{\rule{0.5\linewidth}{0.5pt}}\\[6pt]
		{\small Generated automatically\enspace\textbullet\enspace
			EV Charge DZ Demo Build\enspace\textbullet\enspace April--May~2026}
	\end{center}
	
\end{document}
